\section{SABRE 基线复现实验}

自研路由算法之前，报告先复现 SABRE。这样安排有两个原因：一是 SABRE 本身就是量子
线路映射领域常用的启发式方法，适合作为量子信息基础大作业的经典算法起点；二是后续 NPQR 的
效果需要放在同一批线路、同一类拓扑和同一组指标下讨论。没有稳定基线，单独报告
一个 SWAP 数很难说明算法质量。

\subsection{算法思路}

SABRE 的核心思想是围绕“前沿门”逐步改善映射。在线路依赖关系中，所有前驱已经
完成的门构成当前前沿。若前沿中的双比特门已经落在硬件边上，路由器直接执行该门；
若它的两个逻辑量子比特在当前映射下不相邻，路由器就在相关物理邻边中选择一个
SWAP，更新映射后继续检查前沿。

这个过程把路由问题转化为局部搜索。每一步不尝试枚举完整线路的所有映射序列，而是
用当前前沿门和少量后续门估计候选 SWAP 的价值。设前沿门集合为 $F$，当前映射为
$\pi$，硬件图距离为 $d_G$，一个常用的前沿距离代价可以写成
\[
  H_{\mathrm{front}}(\pi)=
  \sum_{(q_i,q_j)\in F} d_G\bigl(\pi(q_i),\pi(q_j)\bigr).
\]
若候选 SWAP 使这个代价下降，前沿门离可执行状态就更近。SABRE 的 lookahead 和
decay 变体会进一步考虑后续门或重复移动惩罚；basic 变体保留最直接的距离启发式，
适合作为固定、可复现的质量基线。

\subsection{任务定义}

候选 SWAP 不需要从全部物理边中盲目选择。SABRE 通常先找出与前沿门相关的逻辑量子
比特，再把这些逻辑量子比特当前所在物理位置的邻边作为候选集合。这样做减少了搜索
分支，也把每一步选择集中在最可能推进当前线路的位置。

候选动作的评价具有明显的“当前与未来”权衡。只看当前门时，最短路径上的某个 SWAP
往往很有吸引力；但这个 SWAP 可能把另一个即将参与后续门的逻辑量子比特移远。SABRE
的启发式价值在于，它用前沿门和扩展集合共同估计一个动作，而不是只为当前门走一步
最短路。这也是本项目后续设计自研方法时保留“前沿门”“候选动作”“物理距离”这些
概念的原因。

\subsection{复现设置}

本项目复现 SABRE 时采用固定口径。输入线路使用 OpenQASM 2 表示，硬件拓扑使用
Tokyo 风格耦合图；映射预处理先把逻辑量子比特放入目标拓扑，再执行 SABRE 路由。
路由阶段比较 basic、lookahead 和 decay 三种启发式，随机种子固定为 42，单次基线
使用 1 次 trial；输出统一统计完成状态、SWAP 数和线路深度。后续主实验中的
“SABRE basic”即采用这一固定口径。

表 \ref{tab:quick-matrix} 给出小规模复现结果。GHZ5 在三种启发式下均不需要额外
SWAP，说明其交互结构已经适配当前拓扑。QFT5 和 QFT10 的长距离双比特交互更多，
三种启发式之间出现差异；在 QFT10 上，lookahead 的 SWAP 数低于 basic，说明后续门
信息能改善局部选择。报告后续仍选 basic 作为主要对照，是因为它定义清晰、参数少，
更适合作为固定基线；lookahead 和 decay 则用于说明 SABRE 本身也存在启发式设计
空间。

\begin{table}[H]
  \centering
  \caption{SABRE 快速基线矩阵}
  \label{tab:quick-matrix}
  \begin{tabular}{llrrr}
    \toprule
    电路 & 启发式配置 & SWAP 数 & 深度 & 是否完成 \\
    \midrule
    GHZ5 & basic & 0 & 7 & 是 \\
    GHZ5 & lookahead & 0 & 7 & 是 \\
    GHZ5 & decay & 0 & 7 & 是 \\
    QFT5 & basic & 8 & 53 & 是 \\
    QFT5 & lookahead & 7 & 55 & 是 \\
    QFT5 & decay & 7 & 55 & 是 \\
    QFT10 & basic & 42 & 168 & 是 \\
    QFT10 & lookahead & 29 & 156 & 是 \\
    QFT10 & decay & 30 & 167 & 是 \\
    \bottomrule
  \end{tabular}
\end{table}


\subsection{结果分析}

SABRE 复现给后续工作提供了三层参照。第一，它验证了线路解析、拓扑建模、映射更新
和 SWAP 统计这条基本实验链路可以运行。第二，它提供了传统启发式的质量水平，使
NPQR 的 SWAP 数、深度和耗时可以在同一指标下比较。第三，它暴露了经典启发式的
局部性：当线路包含密集长距离交互时，单步启发式很容易受到当前前沿门支配。后续
方法章围绕初始映射、束搜索和模型评分展开，正是为了在保留 SABRE 可解释结构的
同时扩大搜索视野。
